\bookchapter{第二十三章：編程之思與自動化啓蒙}

\chapterintro{前兩章，吾輩爲手機築起了安全之“邊界牆”，亦見識了牆內日益聰明之“AI大腦”。今，吾輩解鎖手機之終極潛能之一：令其如程序員般“思考”，並自動完成君設定之連續任務。君或曰：“吾不會寫代碼，此與吾何干？” 放心，本章所講根本不需君寫一行代碼——現代手機已把“編程”之門檻降至拖拖拽拽即可搞定之程度}

\booksection{一、自動化到底是個啥？}

以大白話言之，自動化即：\textbf{令手機於某條件滿足時，自動爲君做一連串事。}君不需每次手動操作，亦不必記住該做何——手機會自斷時機，自執行。

舉例言之：每日晨君醒來，要手動關掉勿擾模式、打開微信觀消息、打開天氣App觀今日穿何、打開網易雲放首歌。若將此等步驟打包成一個自動化流程，設置成“每日晨7點自動執行”，則君醒後只需拿起手機，一切皆已備好。想想即爽？

自動化本質上乃“若……則……”之邏輯鏈，編程中之if-else。學明此思維方式，縱君一生不寫代碼，處理問題之條理性亦將強甚多。核心要素：

\begin{enumerate}
  \item \textbf{觸發條件}：何時始？可爲"每日晨7點"、"當吾連至家中Wi-Fi"、"當吾至某地"、"當手機電量低於20\%"。
  \item \textbf{執行動作}：具體做何？可爲發一條微信、打開某App、調音量、換壁紙、播音樂、發通知……凡君能於手機上手動做者，自動化皆能做。
  \item \textbf{中間邏輯}（可選然甚強大）：若情況A則做X，否則做Y。譬如"若為工作日則播新聞播客，週末則播音樂"，此曰條件分支。
\end{enumerate}

\booksection{二、iOS“快捷指令”——蘋果用戶之自動化神器}

蘋果之“快捷指令”（Shortcuts）App乃自帶之應用，不需下載。其核心概念乃將各種操作如搭積木般拼起。先言數個關鍵概念：

\begin{itemize}
  \item \textbf{"快捷指令"}：一個完整之自動化流程，可手動點擊運行，亦可經Siri語音喚醒。譬如君曰"嘿Siri，回家模式"，其即自動執行君設定好之一連串操作。
  \item \textbf{"自動化"}：由觸發條件自動運行之指令，不需君手動觸發。譬如"當手機連至車載藍牙時，自動打開地圖App"。自動化於"自動化"標籤頁下設之，與"快捷指令"分開。
  \item \textbf{"動作"}：最小之執行單元。打開App乃一動作，發送消息乃一動作，獲取天氣乃一動作，調整亮度乃一動作。君可將成百上千個動作串於一處。
\end{itemize}

予君觀數個實實在在好用之例，君隨做一遍即明：

\booksubsection{實例一：一鍵“早安流程”}

每日晨點之（或對Siri曰“早上好”），手機將自動：獲取今日天氣朗讀天氣與氣溫顯示今日之日曆安排打開君常用之音樂App始播放。操作步驟：新建快捷指令添加動作“獲取當前天氣”添加“朗讀文本”添加“獲取日曆事件”添加“顯示提醒”添加“打開App”選網易雲音樂。即此簡單，全程拖拽。

\booksubsection{實例二：自動“省電模式”}

於“自動化”標籤頁，新建自動化選“電池電量”設“低於30\%”添加動作“設定低電量模式爲開啓”、“降低屏幕亮度至30\%”、“關閉藍牙”。從此手機電量低於30\%即自動入省電狀態，不需君操心。

\booksubsection{實例三：掃碼自動記賬}

每次君以微信/支付寶掃碼支付後，手機將讀取支付通知，自動將金額與商家記至一備忘錄或Numbers表格中。此需用快捷指令之“自動化”中之“收到短信/通知”觸發條件（iOS 17+支持讀取錢包交易），配“追加至Numbers”動作。對於欲記賬又懶於手動記者，直如救星。

尚有一多人不知之巧用：君可將快捷指令添至主屏幕，變爲一App圖標。譬如做一“打車回家”之快捷指令，圖標放桌面上，點之即徑開滴滴並將目的地設爲家，比君自開App慢慢輸入快多矣。

\booksection{三、安卓自動化——Tasker和它之朋友們}

安卓此邊，情況稍繁，然天花板亦更高。先言簡單者：

\booksubsection{入門級：自動點擊器}

多數安卓手機自帶“自動點擊器”或可於應用商店下載。其理甚簡：君手動操作一遍，其將君所點之座標與時間間隔錄下，而後可高速回放。譬如某App每日要領簽到獎勵，君錄一遍點擊流程，此後每日自動跑一遍即罷。缺點乃：其只認座標，若界面佈局變了（譬如彈了個廣告），其即點歪矣。

\booksubsection{進階級：Tasker——安卓自動化之王}

Tasker乃安卓平臺上公認最強之自動化工具，常年霸榜。實話言其學習曲線稍陡，界面不算好看，然一旦上手，君將覺手機忽變得無所不能。Tasker之核心亦“觸發任務”之式：

\begin{itemize}
  \item \textbf{配置文件（Profile）}：iOS之"到達/離開位置"觸發用者爲地理圍欄技術，準確度取決於GPS信號與系統優化。有時君已到家十分鐘方觸發。建議用Wi-Fi連接或基站觸發代位置觸發，更可靠亦更省電。
  \item \textbf{任務（Task）}：定義"觸發後做何"。一個任務可含多個動作，按序執行。動作類型包括：調整系統設置（音量、亮度、Wi-Fi、藍牙、飛行模式）、發送短信、HTTP請求、執行Shell命令、控制媒體播放、彈出對話框等等。
\end{itemize}

予君數個Tasker之經典實用場景：

\begin{itemize}
  \item \textbf{到家自動開Wi-Fi}：配置觸發條件爲"連接至特定基站"（因君於家附近會連至固定之基站）任務"開啓Wi-Fi"。不用GPS亦不耗電，手機一進小區範圍即自動開Wi-Fi。
  \item \textbf{會議模式}：觸發條件爲"日曆事件標題含'會議'"任務"設置靜音模式、關閉媒體音量、自動回覆微信（需插件配合）"。再也不用憂開會時手機忽響。
  \item \textbf{深夜護眼}：觸發條件爲"時間22:00-06:00"任務"開啓護眼模式、降低亮度至20\%、將屏幕設爲黑白（灰度模式）"。黑白屏幕對刷短視頻之慾望有奇效——畫面一灰，君即不想刷矣。
  \item \textbf{耳機插入自動播客}：觸發條件爲"有線耳機插入/藍牙耳機連接"任務"打開小宇宙App並始播放最新一集"。戴上耳機即能聽，不用解鎖手機覓App。
\end{itemize}

若君覺Tasker上手太難，可試其簡化版兄弟：\textbf{MacroDroid}。界面友好許多，模板豐富，適不願折騰之用戶。另外今多數國產安卓系統（小米之MIUI、華爲之鴻蒙、OPPO之ColorOS）皆自帶“場景”或“智能助手”功能，雖不及Tasker強大，然日常自動化完全夠用。

\booksection{四、自動化思維之日常應用}

學自動化最大之收穫，實則非省那點時間，乃\textbf{改君思考問題之方式}。一旦君習慣“若發生X，即自動做Y”之思維模式，君將覺生活中一大堆重複勞動皆可被省去。譬如：

\begin{itemize}
  \item 每次拍完照片，自動備份至網盤，然後將本地照片壓縮成小尺寸以省空間。
  \item 每日晨自動自日曆讀取今日之安排，生成一段文字經微信發與家人。
  \item 收到快遞簽收短信後，自動於備忘錄中添一條待辦"取快遞"。
  \item 手機電量充滿時自動發語音提醒拔掉充電器（避免過充傷電池）。
  \item 每日下班時間至，若尚於公司附近（經基站定位），自動提醒"該下班矣"。
\end{itemize}

此等聽來甚高級，然實則君只要掌握“觸發+動作”此二基本概念，配合快捷指令或Tasker，皆能做出。關鍵是\textbf{君要先意識到“此事可以自動化”}——多人非不會，乃根本未往那方面想。

\booksection{五、踩坑指南：自動化之常見問題}

搞自動化亦非一帆風順，有幾個坑吾先爲君踩過矣：

\begin{itemize}
  \item \textbf{位置觸發或不靈}：iOS之"到達/離開位置"觸發用之為地理圍欄技術，準確度取決於GPS信號與系統優化。有時君已到家十分鐘方觸發。建議用Wi-Fi連接或基站觸發代位置觸發，更可靠亦更省電。
  \item \textbf{快捷指令之"自動化"需授權}：iOS裏，"自動化"標籤下之觸發多數需君手動確認方能運行（iOS之安全機制）。然Wi-Fi連接、藍牙連接、鬧鐘停止此數個觸發爲可真正全自動運行者，不需確認。
  \item \textbf{Tasker需開啓無障礙權限}：多數高級功能（譬如模擬點擊）需無障礙服務權限，此將令手機於設置中顯一常駐通知。此乃安卓之安全設計，無法去之。
  \item \textbf{勿將自動化搞得太繁}：剛開始做自動化，從"2-3個動作"之小流程始。勿一上來即搞一個20步之超級流程——調試起來將瘋。
\end{itemize}

\textbf{小貼士：自動化入門之正確姿勢}\par

先費5分鐘思之：君每日於手機上重複最多之3件事爲何？譬如每日晨之例行操作、每次連藍牙耳機後之習慣、每次出門前之設置……自此等高頻小事入手，各做一簡單之自動化。哪怕每個自動化只省君5秒鐘，日積月累亦可觀。且做出後之滿足感，真會上癮。

\begin{tcbpractice}{練一練：三個入門自動化挑戰}
選一個最適合君者完成，做完君即入門矣：

\begin{enumerate}
  \item \textbf{挑戰一：早安自動化}：iOS：於快捷指令App中新建快捷指令添加"獲取當前天氣""朗讀文本""打開App選網易雲"。完成後君可對Siri曰"早上好"以觸發。安卓：於Tasker中新建任務添加"獲取天氣""朗讀""打開App"設觸發時間爲每日7:00。
  \item \textbf{挑戰二：低電量自動省電}挑戰二：低電量自動省電
  \item \textbf{挑戰三：連接車載藍牙自動導航}：iOS：自動化新建"藍牙"選擇君之車載藍牙添加"打開App選高德地圖"+"媒體音量調至80\%"。安卓：TaskerProfile狀態藍牙已連接選擇車載藍牙任務"打開高德"+"媒體音量調高"。
\end{enumerate}

\end{tcbpractice}

\begin{tcbinfo}{案例 3：自動化"早安流程"}
\begin{itemize}
  \item \textbf{場景}：每日 7:00 自動播報天氣 + 關閉靜音 + 打開新聞。
  \item \textbf{步驟}：iOS: 快捷指令自動化時間觸發添加"獲取天氣"+"設定音量"+"打開 App";Android: MacroDroid觸發器(時間 7:00)動作(媒體音量/打開 URL/TTS 朗讀天氣)。
  \item \textbf{預期果}：免安裝即可實現，降低上手門檻。
\end{itemize}

\end{tcbinfo}
